% LTeX: language=de

\subsection{Grundlagen}

\vspace{-0.5cm}

\raisebox{-\height}{\incldrawio{01_Umkehrfunktion_Veranschaulichung.drawio}}

\begin{Definition}{Umkehrfunktion}{}
    Die \textbf{Umkehrfunktion $f^{-1}$} ($\neq \dfrac{1}{f}$) einer Funktion $f$ ordnet jedem Element der Wertemenge $W_{f}$ das zugehörige Element aus $D_{f}$ zu.
\end{Definition}

\raisebox{-\height}{
    \begin{tikzpicture}
        % Set variables
        \pgfmathsetmacro{\axisPosWidth}{7};
        \pgfmathsetmacro{\axisNegWidth}{-7};
        \pgfmathsetmacro{\axisPosHeight}{4};
        \pgfmathsetmacro{\axisNegHeight}{-3};
        \pgfmathsetmacro{\xIncrement}{1.0};
        \pgfmathsetmacro{\yIncrement}{1.0};
        \pgfmathsetmacro{\xScale}{1};
        \pgfmathsetmacro{\yScale}{2};
        \pgfmathsetmacro{\gridxIncrement}{0.5};
        \pgfmathsetmacro{\gridyIncrement}{0.5};

        % Axis/grid limits
        \pgfmathsetmacro{\xUpperLimitAxis}{\xIncrement * floor(\axisPosWidth / \xIncrement)}
        \pgfmathsetmacro{\xLowerLimitAxis}{\xIncrement * ceil(\axisNegWidth / \xIncrement)}
        \pgfmathsetmacro{\yUpperLimitAxis}{\yIncrement * floor(\axisPosHeight /\yIncrement)}
        \pgfmathsetmacro{\yLowerLimitAxis}{\yIncrement * ceil(\axisNegHeight / \yIncrement)}

        \pgfmathsetmacro{\xUpperLimitGrid}{\gridxIncrement * floor(\axisPosWidth / \gridxIncrement)}
        \pgfmathsetmacro{\xLowerLimitGrid}{\gridxIncrement * ceil(\axisNegWidth / \gridxIncrement)}
        \pgfmathsetmacro{\yUpperLimitGrid}{\gridyIncrement * floor(\axisPosHeight / \gridyIncrement)}
        \pgfmathsetmacro{\yLowerLimitGrid}{\gridyIncrement * ceil(\axisNegHeight / \gridyIncrement)}

        % Axis environment (PGFPlots)
        \begin{axis}[
            xmin=\xLowerLimitGrid, xmax=\xUpperLimitGrid,
            ymin=\yLowerLimitGrid, ymax=\yUpperLimitGrid,
            axis lines=middle,
            axis line style={->, >=stealth},
            grid=both,
            xlabel={$x$},
            ylabel={$y$},
            xlabel style={at={(axis description cs:1,0.55)}, anchor=north},
            ylabel style={at={(axis description cs:0.55,1)}, anchor=south},
            xtick={\xLowerLimitAxis,\xLowerLimitAxis+1,...,\xUpperLimitAxis},
            ytick={\yLowerLimitAxis,\yLowerLimitAxis+1,...,\yUpperLimitAxis},
            xticklabel={\pgfmathparse{\tick*\xScale}\pgfmathprintnumber{\pgfmathresult}},
            yticklabel={\pgfmathparse{\tick*\yScale}\pgfmathprintnumber{\pgfmathresult}},
            tick label style={font=\scriptsize},
            x=1cm,
            y=1cm,
            minor tick num=1,
            scale only axis,
            grid=both,
            restrict y to domain=\yLowerLimitGrid:\yUpperLimitGrid,
        ]

        % G_f
        \addplot[thick, myb, samples=1000, domain=\xLowerLimitGrid:0.96*\xUpperLimitGrid]
            {(0.5*\x*\xScale + 1)*(1/\yScale)};
        \node[anchor=north] at (axis cs:{5*(1/\xScale)},{3*(1/\yScale)}) {\textcolor{myb}{$G_f$}};

        % G_f^{-1}
        \addplot[thick, myr, samples=1000, domain=\xLowerLimitGrid:0.96*\xUpperLimitGrid]
            {(2.0*\x*\xScale - 2)*(1/\yScale)};
        \node[anchor=north] at (axis cs:{3*(1/\xScale)},{6*(1/\yScale)}) {\textcolor{myr}{$G_f^{-1}$}};

        % Winkelhalbierende
        \addplot[thick, myg, samples=1000, domain=\xLowerLimitGrid:0.96*\xUpperLimitGrid]
            {(\x*\xScale)*(1/\yScale)};

        % Qudranten
        \node[anchor=north] at (axis cs:{6.5*(1/\xScale)},{8*(1/\yScale)}) {\textcolor{black}{$\Rom{1}$}};
        \node[anchor=north] at (axis cs:{-6.5*(1/\xScale)},{8*(1/\yScale)}) {\textcolor{black}{$\Rom{2}$}};
        \node[anchor=north] at (axis cs:{-6.5*(1/\xScale)},{-5*(1/\yScale)}) {\textcolor{black}{$\Rom{3}$}};
        \node[anchor=north] at (axis cs:{6.5*(1/\xScale)},{-5*(1/\yScale)}) {\textcolor{black}{$\Rom{4}$}};

        \end{axis}
    \end{tikzpicture}
}\\[3.0ex]

\begin{itemize}
    \item $G_f^{-1}$ entsteht aus $G_f$ durch Spiegelung an der \texthl{\textbf{Winkelhalbierenden}} \\ des \Rom{1} und \Rom{3} Qudranten.
    \item $D = W_{f}$ \qquad und \qquad $W_{f^{-1}} = D_{f}$
    \item \begin{minipage}[t]{0.65\linewidth}
        \raisebox{-\height}{
            \begin{tikzpicture}
                % Set variables
                \pgfmathsetmacro{\axisPosWidth}{5};
                \pgfmathsetmacro{\axisNegWidth}{-3};
                \pgfmathsetmacro{\axisPosHeight}{3};
                \pgfmathsetmacro{\axisNegHeight}{-2.5};
                \pgfmathsetmacro{\xIncrement}{1.0};
                \pgfmathsetmacro{\yIncrement}{1.0};
                \pgfmathsetmacro{\xScale}{1};
                \pgfmathsetmacro{\yScale}{1};
                \pgfmathsetmacro{\gridxIncrement}{0.5};
                \pgfmathsetmacro{\gridyIncrement}{0.5};

                % Axis/grid limits
                \pgfmathsetmacro{\xUpperLimitAxis}{\xIncrement * floor(\axisPosWidth / \xIncrement)}
                \pgfmathsetmacro{\xLowerLimitAxis}{\xIncrement * ceil(\axisNegWidth / \xIncrement)}
                \pgfmathsetmacro{\yUpperLimitAxis}{\yIncrement * floor(\axisPosHeight /\yIncrement)}
                \pgfmathsetmacro{\yLowerLimitAxis}{\yIncrement * ceil(\axisNegHeight / \yIncrement)}

                \pgfmathsetmacro{\xUpperLimitGrid}{\gridxIncrement * floor(\axisPosWidth / \gridxIncrement)}
                \pgfmathsetmacro{\xLowerLimitGrid}{\gridxIncrement * ceil(\axisNegWidth / \gridxIncrement)}
                \pgfmathsetmacro{\yUpperLimitGrid}{\gridyIncrement * floor(\axisPosHeight / \gridyIncrement)}
                \pgfmathsetmacro{\yLowerLimitGrid}{\gridyIncrement * ceil(\axisNegHeight / \gridyIncrement)}

                % Axis environment (PGFPlots)
                \begin{axis}[
                    xmin=\xLowerLimitGrid, xmax=\xUpperLimitGrid,
                    ymin=\yLowerLimitGrid, ymax=\yUpperLimitGrid,
                    axis lines=middle,
                    axis line style={->, >=stealth},
                    grid=both,
                    xlabel={$x$},
                    ylabel={$y$},
                    xlabel style={at={(axis description cs:1,0.55)}, anchor=north},
                    ylabel style={at={(axis description cs:0.45,1)}, anchor=south},
                    xtick={\xLowerLimitAxis,\xLowerLimitAxis+1,...,\xUpperLimitAxis},
                    ytick={\yLowerLimitAxis,\yLowerLimitAxis+1,...,\yUpperLimitAxis},
                    xticklabel={\pgfmathparse{\tick*\xScale}\pgfmathprintnumber{\pgfmathresult}},
                    yticklabel={\pgfmathparse{\tick*\yScale}\pgfmathprintnumber{\pgfmathresult}},
                    tick label style={font=\scriptsize},
                    x=1cm,
                    y=1cm,
                    minor tick num=1,
                    scale only axis,
                    grid=both,
                    restrict y to domain=\yLowerLimitGrid:\yUpperLimitGrid,
                ]

                % G_g
                \addplot[thick, myb, samples=1000, domain=\xLowerLimitGrid:0.96*\xUpperLimitGrid]
                    {((\x*\xScale)^2)*(1/\yScale)};
                \node[anchor=north] at (axis cs:{2*(1/\xScale)},{3*(1/\yScale)}) {\textcolor{myb}{$G_g$}};

                % G_g^{-1}
                \addplot[thick, myr, samples=1000, domain=0:0.96*\xUpperLimitGrid]
                    {(sqrt(\x*\xScale))*(1/\yScale)};
                \addplot[thick, myr, samples=1000, domain=0:0.96*\xUpperLimitGrid]
                    {(-sqrt(\x*\xScale))*(1/\yScale)};
                \node[anchor=north] at (axis cs:{3*(1/\xScale)},{1.5*(1/\yScale)}) {\textcolor{myr}{$G_g^{-1}$}};

                \end{axis}
            \end{tikzpicture}
        }\\[3.0ex]
        \end{minipage}
        \begin{minipage}[t]{0.3\linewidth}
            \vspace{2.5cm}
            $g$ ist auf $\rbrack \ -\infty; \ 0 \ \rbrack$ umkehrbar sowie auf $\lbrack \ 0; \ +\infty \ \lbrack$
        \end{minipage}\\
        Eine Funktion ist genau dann umkehrbar, falls jdes Element der Wertemenge \textbf{genau einmal} getroffen wird. \\
        \vspace{-1.0cm}
        \begin{tightcenter}
            $\Leftrightarrow$
        \end{tightcenter}
        \vspace{-0.5cm}
        Eine Funktion ist genau dann umkehrbar, falls sie \texthl{\textbf{streng monoton}} (zunehmend oder abhnehmend) ist.
    \item es gilt: \begin{minipage}[t]{0.80\linewidth}
        \begin{flalign*}
            f(f^{-1}(x)) &= x \\[1.0ex]
            f^{-1}(f(x)) &= x
        \end{flalign*}
    \end{minipage}
\end{itemize}